\HuazhongGuide{本表参照以下提纲撰写，要求填写内容真实准确。}

\section{\textbf{工作基础}：与本项目相关的研究工作积累和已取得的研究工作成绩。}

\subsection{申请人简历与联系方式}
\HZField{applicant}，\HZField{sex}，\HZField{ethnicity}，出生年月\HZField{birth}，\HZField{degree}，\HZField{professional-title}，\HZField{supervisor}。工作单位：\HZField{organization}；专业：\HZField{major}。联系电话：\HZField{phone}；电子邮箱：\HZField{email}。通讯地址：\HZField{address}，邮编\HZField{postcode}。

\textbf{教育经历。}\HZField{education}

\textbf{工作经历。}\HZField{employment}

% 迁移来源：extraTex/3.1.研究基础.tex；与本项目相关的研究工作积累
\subsection{与本项目相关的研究工作积累}

申请人\HZField{applicant}为\HZField{school}\HZField{professional-title}、\HZField{supervisor}，具备机器学习、表示学习、问答及软件系统研发积累。本课题所需的知识证据组织、检索结果分析与可复现实验，与申请人已有工作存在方法联系。前期研究主要支撑表征、检索及实验实现能力，教育情境标注和智能体风险监测仍是本项目拟开展的研究。

申请人主持的“面向政企数字化应用的智能知识库与问答系统研发”横向课题，提供了知识库建设任务的工程衔接。该课题依合同构建知识库，其客户场景与本项目的教育风险评价对象不同；项目合同中的资料与成果不能未经授权直接作为开放评测数据。教育任务与伦理风险标签将在本课题中独立设计和核验。

\subsection{研究条件与实施可行性}\label{subsec:conditions}
申请人可使用A40显卡8张、RTX 3090 Ti显卡6张及Linux、Python、深度学习、Git、LaTeX工具。以已有模型推理、重放和评价为主，不训练基础模型、不新购大型设备、不假定算力长期独占。问答、检索及知识库积累支撑实施，尚无可复用的本项目轨迹、脚本或先导效果证据。项目拟组织具有相关教学经验的人员复核任务与评价结果，先开展公开资料和模拟环境实验；涉及真实课程及师生数据的验证，在取得相应授权并完成适用的伦理审查程序后开展。模拟案例不用于推断真实学习效果。拟开展基地合作，专属设备与教学数据的使用需另行取得授权；真人参与须与成绩分离，退出不影响权益。

\textbf{组织实施安排。}申请人负责总体设计、关键方法、质量复核和成果组织。主要成员为孔哲、张建宏、杨铄、孟星宇、桑乐、吕中原和于得武，均为新疆大学在读硕士研究生。孔哲拟负责公开课程资料、许可核对与任务构建；张建宏拟负责风险规则、标注手册及质量核查；杨铄拟负责轨迹采集与对照复现；孟星宇拟负责跨任务评测及稳健性分析；桑乐拟负责情境约束消融与监测校准；吕中原拟负责分级干预及风险—效用分析；于得武拟负责代码集成、实验日志及统计复核。

成员按任务家族交叉复核，同一任务的生成者不承担其独立标注，仲裁者不兼任该任务的初始标注者。正式标注前开展手册培训与一致性试标；教学规则由具备课程理解能力的人员复核，相关外部人选和实验室合作分工在研究启动阶段协商确定。研究助理费用按七名成员分阶段承担工作的累计16人月测算，外聘专家咨询按4人天测算，实际核算按任务记录及单位规定执行。

\subsection{研究风险与应对措施}\label{subsec:risks}
歧义集中时补充正反例、手册和敏感性分析，必要时收缩为授权明确任务；教学复核人员未落实则限为合成协议验证。增益不稳定时检查信息量、长度和构造捷径，不换测试集追求有利结论。干预负担过高时收缩自动覆盖，保留确定性权限保护，争议转人工且不删除失败任务。授权或模型成本受限时使用公开材料及可本地复现模型，记录版本费用，减少公平性、双智能体扩展，保证一年内主比较。本次按一般课题申请5万元，直接费用4万元、间接费用1万元，优先保障标注复核、核心对照与案例研究。立项后按指南首拨60\%、中期评价通过后拨付余款，前期以现有设备完成准备和预实验；未获后续拨款前不扩展次要任务。

\section{\textbf{申请人代表作}：近期已获得与本项目相关的科研成果。须注明作者、成果名称、形式（如论文、专著、研究报告等）、发表刊物或出版社名称、发表或出版时间等信息。奖励情况也须详细列出全部受奖人员、奖励名称等级、授奖年等。\textbf{（限10 项）}}

% 迁移来源：extraTex/3.1.研究基础.tex；已取得的研究工作成绩
\subsection{已取得的研究工作成绩}

Feng Yan, Zhe Li, Wushour Silamu, Yanbing Li. Knowledge-aware image understanding with multi-level visual representation enhancement for visual question answering发表于\textit{Machine Learning}，正式卷期为2024年113卷6期3789—3805页\cite{yan2024knowledge}。申请人为第一作者、通讯作者，DOI：\url{https://doi.org/10.1007/s10994-023-06426-6}。该工作涉及知识与多层表示的结合，可为本课题组织证据和分析问答行为提供方法经验，教育智能体可信评价的有效性将在本项目中验证。

Lihan Tang, Liejun Wang, Gang Wang, Mengyuan Sun, Feng Yan. Feature Fusion Mamba Hashing via Decoupling for Cross-Modal Retrieval发表于\textit{IEEE Signal Processing Letters}，2026年33卷609—613页\cite{tang2026hashing}。申请人为通讯作者，DOI：\url{https://doi.org/10.1109/LSP.2026.3652952}。跨模态检索积累有助于本课题的证据组织和检索分析，但不将检索精度作为教育效果指标。

Pengfei Xu, Feng Yan, Yuqing Zhou, Yuancheng Liu, Chuanrui Wu. RSHR+: Progressive Question-Conditioned Visual Calibration and Structured State-Space Reasoning for Remote Sensing Visual Question Answering，\textit{Expert Systems with Applications}，2026年328卷，文章号132953，DOI：\url{https://doi.org/10.1016/j.eswa.2026.132953}。该期刊论文已有在线题录，正式卷期日期为2026年10月；申请人为通讯作者。问题条件化表示与长序列推理的经验，可为本课题处理任务约束、轨迹上下文及对照分析提供方法基础，教育风险监测成效需在本课题中另行检验。

申请人在上述三项论文中负责研究构思、方法设计、实验指导、结果分析及论文撰写与修改。

本课题尚无可直接复用的教育轨迹数据、风险评价脚本和消融实验日志。EduAgentTrust评测集和监测原型为本次研究的预期产出，其性能与适用范围将在研究过程中检验。

\section{\textbf{承担项目}：申请人承担的各级各类科研项目情况。包括项目名称和编号、资助机构、资助金额、起止年月、与本项目的关系等。\textbf{（限10 项）}}

% 迁移来源：extraTex/3.3.承担项目情况.tex；全文
\subsection{相关在研项目}
申请人目前承担的项目中，与本课题知识证据组织和系统实现直接相关的是“面向政企数字化应用的智能知识库与问答系统研发”，委托单位为新疆数字云创有限公司，横向课题编号202604140013，经费12万元，执行期2026年6月至10月。申请人为负责人，合同任务为构建一个知识库；成果及知识产权由合同双方共同享有。本课题面向高等教育智能体的可信风险表征与监测，不承担该合同的客户系统交付，数据与成果使用分别遵循约定。

申请人主持“多模态融合技术研究及其在新疆产业集群中的应用”，属自治区天池英才引进计划青年博士人才项目，经费50万元，执行期2025年2月至2028年2月，以任务书为依据，无独立项目批准号。申请人负责总体方案、关键模型设计及成果组织。其研究对象为多模态技术与产业应用，可在表示和实验经验上形成支撑；本课题独立研究教育情境约束及轨迹风险，不将产业项目的研究任务重复列为本课题成果。

申请人参与CAAI—蚂蚁科研基金项目“可学习乘积量化与状态空间模型协同的长上下文高效注意力机制研究”，编号202604140006，项目总经费15万元，执行期2026年2月至2027年2月，本人负责高效注意力机制研究。该项目为长上下文处理与运行成本分析提供方法经验；本课题使用已有模型检验教育约束及动作前风险，不重复承担该基金的注意力结构研发任务。

申请人主持自治区高校基本科研业务费科研项目培育类“融合多源数据大模型驱动的光伏发电功率预测方法研究”，编号XJEDU2026P005，经费5万元，其中直接费用3.9万元、间接费用1.1万元，执行期2026年1月至12月。申请人负责总体方案、模型研究及成果组织。该项目提供模型实验和数据质量分析经验，研究对象与评价终点是光伏功率预测；本课题单独建设教育任务、标注及轨迹日志，分别核算经费与成果，不共用受限业务数据。

七名主要参与者的组成及分工见成员表；申请人将逐项核实其与本课题相关的在研项目，依据具体研究任务明确各项目之间的任务、经费及成果边界。

% 迁移来源：extraTex/3.4.完成国基项目情况.tex；原稿中未主持完成国基项目的说明
\subsection{既往主持国基项目情况}

申请人未主持完成国家自然科学基金项目。
